% Part: computability
% Chapter: computability-theory
% Section: def-functions-self-reference

\documentclass[../../../include/open-logic-section]{subfiles}

\begin{document}

\olfileid[tr]{cmp}{thy}{slf}
\olsection{Fonksiyonların Kendine Göndermeyle Tanımlanması}

Fonksiyonları kendileri cinsinden tanımlayabilmek genel olarak yararlıdır.
Örneğin hesaplanabilir $k$, $l$ ve $m$ fonksiyonları verildiğinde sabit nokta
lemması, her $y$ için
\[
  f(y)\simeq\begin{cases}
    k(y) & \text{$l(y)=0$ ise}\\
    f(m(y)) & \text{aksi durumda}
  \end{cases}
\]
denklemini sağlayan kısmi hesaplanabilir bir $f$ fonksiyonunun bulunduğunu
söyler. Daha açık olarak $f$,
\[
  g(x,y)\simeq\begin{cases}
    k(y) & \text{$l(y)=0$ ise}\\
    \cfind{x}(m(y)) & \text{aksi durumda}
  \end{cases}
\]
fonksiyonu alınarak ve ardından sabit nokta lemmasıyla
$\cfind{e}(y)=g(e,y)$ olacak bir $e$ indisi bulunarak elde edilir.

Somut bir örnek olarak ``en büyük ortak bölen'' fonksiyonu
$\fn{gcd}(u,v)$
\[
  \fn{gcd}(u,v)\simeq\begin{cases}
    v & \text{$u=0$ ise}\\
    \fn{gcd}(\fn{mod}(v,u),u) & \text{aksi durumda}
  \end{cases}
\]
ile tanımlanabilir; burada $\fn{mod}(v,u)$, $v$'nin $u$'ya bölümünden kalanı
gösterir. Sabit nokta lemmasına başvurmak $\fn{gcd}$ fonksiyonunun kısmi
hesaplanabilir olduğunu gösterir. Aslında $y$'nin $\tuple{u,v}$ çiftini
kodlaması sağlanarak bu yukarıdaki biçime getirilebilir. Ardından $u$
üzerinde bir tümevarım, $\fn{gcd}$ fonksiyonunun gerçekte tümel olduğunu
gösterir.

Elbette kendine göndermeli tanımların, az önce tartışılan örneklerden çok daha
gösterişli olanları kurulabilir. Programlama dillerinin çoğu fonksiyonların
bir biçimde kendileri cinsinden tanımlanmasını destekler. Bunun, sabit nokta
teoreminin tam genelliği içinde yapmamızı sağladığı gibi bir fonksiyonu,
fonksiyonu hesaplayan algoritmanın bir \emph{indisi} cinsinden tanımlamaktan
biraz daha az çarpıcı olduğuna dikkat edin.

\end{document}

